%% kapa: exam=1 round=34 year=2023 period=1 form=A subject=civil_law subject_ko=민법 seq=39 number=39
%% source: Q-Net 감정평가사 기출문제 (cst00309) artlSeq=5211789 file=raw/1008/5211789/2023년도 제34회 감정평가사 1차시험 1교시.pdf
%% conversion: kordoc 4.24.51 -> tools/gen1.py (auto); review=proofread-sonnet; flags=none
\begin{question}{39}
乙명의의 저당권이 설정되어 있는 甲소유의 X토지위에 Y건물이 신축된 후, 乙의 저당권이 실행된 경우에 관한 설명으로 옳은 것을 모두 고른 것은? (다툼이 있으면 판례에 따름)
\begin{framedtext}
ㄱ. 甲이 Y건물을 신축한 경우, 乙은 Y건물에 대한 경매도 함께 신청할 수 있으나 Y건물의 경매대가에서 우선변제를 받을 수는 없다. \par
ㄴ. Y건물을 甲이 건축하였으나 경매 당시 제3자 소유로 된 경우, 乙은 Y건물에 대한 경매도 함께 신청할 수 있다. \par
ㄷ. Y건물이 X토지의 지상권자인 丙에 의해 건축되었다가 甲이 Y건물의 소유권을 취득하였다면 乙은 Y건물에 대한 경매도 함께 신청할 수 있다. \par
\end{framedtext}
\begin{choices}
  \item ㄴ
  \item ㄱ, ㄴ
  \item ㄱ, ㄷ
  \item ㄴ, ㄷ
  \item ㄱ, ㄴ, ㄷ
\end{choices}
\end{question}
\answer{3}
